\begin{exercisechunk}
\item \textbf{Two derivations of the optimizer identities.}
\label[exercise]{ex:l9-push-through}
\exerciseprofile{2}{2}{2}{\exproof\quad\excalculation}{%
Derive the optimizer from its optimality equations and verify it with a
push-through identity.}
Let $U\in\R^{n\times m}$ and $V\in\R^{n\times D}$, where $V$ has full row
rank, and define
\[
 H=D^{-1}VV^\top,
 \qquad K=H+UU^\top,
 \qquad \alpha=K^{-1}\mathbf Y.
\]
Assume that $H$ is positive definite. Consider
\[
 \min_{\beta\in\R^m,\,b\in\R^D}
 \left\{\frac12\|\beta\|^2+\frac D2\|b\|^2:
 U\beta+Vb=\mathbf Y\right\}.
\]

\begin{enumerate}[(a),beginpenalty=10000]
\item Introduce a Lagrange multiplier in $\R^n$ and use the optimality
equations to prove that the optimizer satisfies
\[
 \widehat\beta=U^\top\alpha,
 \qquad
 \widehat b=D^{-1}V^\top\alpha.
\]

\item Let $A\in\R^{p\times q}$ and $B\in\R^{q\times p}$, and suppose the
two inverses below exist. Starting from
\[
 (I_p+AB)A=A(I_q+BA),
\]
prove the push-through identity
\[
 (I_p+AB)^{-1}A=A(I_q+BA)^{-1}.
\]
Then apply it to prove
\[
 (I_m+U^\top H^{-1}U)^{-1}U^\top H^{-1}
 =U^\top(H+UU^\top)^{-1}.
\]
Conclude that the formula for $\widehat\beta$ in
\cref{eq:beta-hat} equals $U^\top\alpha$.

\item Starting from the formula for $b_\beta$ in \cref{eq:bbeta}, verify
directly that
\[
 \widehat b=D^{-1}V^\top\alpha.
\]
Your calculation should use $VV^\top=DH$ and
$I_n-UU^\top K^{-1}=HK^{-1}$.
\end{enumerate}

\begin{remark}[Why ``push-through''?]
In $(I_p+AB)^{-1}A=A(I_q+BA)^{-1}$, the factor $A$ moves through the
inverse. The matrix being inverted changes from a $p\times p$ matrix to a
$q\times q$ matrix. This identity is especially useful when one of these
two matrices is substantially smaller than the other.
\end{remark}
\end{exercisechunk}
